%% ma: L12-B08-0143
%% lop: 12
%% chuong: 2
%% bai: 8
%% dang: 12.8.48
%% loai: TN4
%% muc: VD
%% dap_an: "B"
%% nguon: "CĐ2.3.pdf giáo viên gửi 27/10/2026 (M07, Lê Minh Khôi – Nguyễn Anh Khôi, Chương II Chủ đề 3 Biểu thức tọa độ của các phép toán vectơ), Câu 59"
%% kien_thuc: "Bài 8 (điều kiện vuông góc, ba điểm thẳng hàng)"
%% trang_thai: cho-duyet
%% ly_do: "Dạng mới 12.8.48 đề xuất, chờ giáo viên duyệt; Đề gốc không có phương án đúng; đã sửa phương án B từ 160 thành 158"
%% ghi_chu: "Đề gốc không có phương án đúng (phương án B ghi L = 160, giá trị đúng là 158); đã sửa phương án B thành 158"
\begin{ex}
Trong không gian $Oxyz$, cho ba điểm $M(0;2;7)$, $T(5;1;-7)$, $P(11;3;6)$. Gọi $d$ là đường thẳng đi qua hai điểm $M$, $P$. Biết điểm $T'(m;n;p)$ là hình chiếu của $T$ lên $d$. Tính giá trị của biểu thức $L=12m+3n+12p$.
\choice{$L=80$}{\True $L=158$}{$L=240$}{$L=320$}
\loigiai{$\overrightarrow{MP}=(11;1;-1)$, $MP^2=123$, $\overrightarrow{MT}=(5;-1;-14)$, $t=\dfrac{68}{123}$.

$T'\left(\dfrac{748}{123};\dfrac{314}{123};\dfrac{793}{123}\right)$, $L=\dfrac{8976+942+9516}{123}=158$.}
\end{ex}
